\documentclass[12pt,a4paper]{article}

% ==================================================
% PAQUETES
% ==================================================
\usepackage[spanish]{babel}
\usepackage[T1]{fontenc}
\usepackage{fontspec}
\setmainfont{Times New Roman}

\usepackage{graphicx}
\usepackage{geometry}
\usepackage{setspace}
\usepackage{fancyhdr}
\usepackage{eso-pic}
\usepackage{titlesec}
\usepackage{enumitem}
\usepackage{longtable}
\usepackage{array}
\usepackage{hyperref}
\usepackage{hanging}

% ==================================================
% CONFIGURACIÓN APA 7
% ==================================================
\geometry{
	left=2.54cm,
	right=2.54cm,
	top=2.54cm,
	bottom=2.54cm
}

\doublespacing

\setlength{\parindent}{1.27cm}
\setlength{\parskip}{0pt}
\setlength{\headheight}{15pt}

% ==================================================
% NUMERACIÓN APA
% ==================================================
\pagestyle{fancy}
\fancyhf{}
\fancyhead[R]{\thepage}
\renewcommand{\headrulewidth}{0pt}

% ==================================================
% CONFIGURACIÓN DE TÍTULOS APA
% ==================================================
\titleformat{\section}
{\normalfont\bfseries\centering}
{}
{0pt}
{}

\titleformat{\subsection}
{\normalfont\bfseries}
{}
{0pt}
{}

\titlespacing*{\section}{0pt}{0pt}{12pt}
\titlespacing*{\subsection}{0pt}{12pt}{6pt}

\setlist[itemize]{
	noitemsep,
	topsep=0pt,
	leftmargin=1.27cm
}

% ==================================================
% FONDOS PDF
% ==================================================
\newcommand{\FondoPortada}{
	\AddToShipoutPictureBG*{
		\AtPageLowerLeft{
			\includegraphics[
			width=\paperwidth,
			height=\paperheight
			]{plantillasTrabajosUNEMI/portadaLibroUNEMI.pdf}
		}
	}
}

\newcommand{\FondoDocumento}{
	\AddToShipoutPictureBG{
		\AtPageLowerLeft{
			\includegraphics[
			width=\paperwidth,
			height=\paperheight
			]{plantillasTrabajosUNEMI/plantillaA4PaginaDocumento.pdf}
		}
	}
}

\newcommand{\FondoFinal}{
	\AddToShipoutPictureBG*{
		\AtPageLowerLeft{
			\includegraphics[
			width=\paperwidth,
			height=\paperheight
			]{plantillasTrabajosUNEMI/plantillaA4PaginaFinal.pdf}
		}
	}
}

% ==================================================
% INICIO DOCUMENTO
% ==================================================
\begin{document}
	
	% ==================================================
	% PORTADA LIMPIA
	% ==================================================
	\FondoPortada
	
	\begin{titlepage}
		\thispagestyle{empty}
		\mbox{}
	\end{titlepage}
	
	\ClearShipoutPictureBG
	
	% ==================================================
	% PORTADA APA 7
	% ==================================================
	\FondoDocumento
	
	\setcounter{page}{2}
	
	\newpage
	
	\thispagestyle{fancy}
	
	\vspace*{3cm}
	
	\begin{center}
		
		\textbf{Historia de las TI}\\[2cm]
		
		Ángel Fernando Calero Maldonado\\[0.5cm]
		
		Universidad Estatal de Milagro\\[0.5cm]
		
		Facultad de Ciencias e Ingeniería\\[0.5cm]
		
		Carrera de Tecnologías de la Información\\[0.5cm]
		
		Fundamentos de tecnologías de información\\[0.5cm]
		
		Díaz Sandoya Eduardo Luis\\[0.5cm]
		
		23 de abril de 2026
		
	\end{center}
	
	% ==================================================
	% INTRODUCCIÓN
	% ==================================================
	\newpage
	
	\section*{Introducción}
	
	El presente documento académico aborda temas relacionados con el software, su clasificación, la calidad del software y los lenguajes de programación de alto y bajo nivel. Estos contenidos permiten comprender la importancia del software dentro del campo de las tecnologías de la información.
	
	La investigación se desarrolla a partir de la información proporcionada en la imagen base, donde se establecen tres ejes principales: clasificación del software, calidad del software y lenguajes de alto y bajo nivel.
	
	% ==================================================
	% DESARROLLO
	% ==================================================
	\section{Software: Clasificación}
	
	El software puede clasificarse según su función principal dentro de un sistema informático. Entre las categorías más importantes se encuentran el software de sistema, el software de aplicación y el software de programación.
	
	\subsection{Software de sistema}
	
	El software de sistema permite el funcionamiento general del computador y administra los recursos del hardware.
	
	\begin{itemize}
		\item Windows.
		\item Linux.
		\item macOS.
	\end{itemize}
	
	\subsection{Software de aplicación}
	
	El software de aplicación está diseñado para que el usuario realice tareas específicas.
	
	\begin{itemize}
		\item Microsoft Word.
		\item WhatsApp.
		\item Google Drive.
	\end{itemize}
	
	\subsection{Software de programación}
	
	El software de programación permite crear, editar y ejecutar programas informáticos.
	
	\begin{itemize}
		\item Python.
		\item Visual Studio Code.
		\item Eclipse.
	\end{itemize}
	
	\subsection{Tabla comparativa}
	
	\begin{longtable}{|p{4cm}|p{9cm}|}
		\hline
		\textbf{Categoría} & \textbf{Ejemplos} \\
		\hline
		Software de sistema & Windows, Linux, macOS. \\
		\hline
		Software de aplicación & Microsoft Word, WhatsApp, Google Drive. \\
		\hline
		Software de programación & Python, Visual Studio Code, Eclipse. \\
		\hline
	\end{longtable}
	
	\subsection{Software según licencia}
	
	El software también puede clasificarse según el tipo de licencia.
	
	\begin{itemize}
		\item \textbf{Software libre:} permite usar, estudiar, modificar y distribuir el programa.
		\item \textbf{Software propietario:} posee restricciones de uso, copia o modificación.
		\item \textbf{Software de código abierto:} permite acceder al código fuente.
		\item \textbf{Software en la nube:} funciona mediante servicios disponibles a través de internet.
	\end{itemize}
	
	% ==================================================
	% CALIDAD DEL SOFTWARE
	% ==================================================
	\section{Calidad del software}
	
	La calidad del software se relaciona con el grado en que un programa cumple con las necesidades del usuario y funciona de manera adecuada.
	
	Según la información base, se deben considerar atributos como funcionalidad, fiabilidad, usabilidad, eficiencia, mantenibilidad y portabilidad.
	
	\subsection{Atributos de calidad}
	
	\begin{itemize}
		\item \textbf{Funcionalidad:} capacidad del software para cumplir las funciones requeridas.
		\item \textbf{Fiabilidad:} capacidad de funcionar correctamente durante un tiempo determinado.
		\item \textbf{Usabilidad:} facilidad con la que el usuario puede utilizar el sistema.
		\item \textbf{Eficiencia:} uso adecuado de recursos como memoria, tiempo y procesamiento.
		\item \textbf{Mantenibilidad:} facilidad para corregir, mejorar o actualizar el software.
		\item \textbf{Portabilidad:} capacidad de funcionar en distintos entornos o plataformas.
	\end{itemize}
	
	\subsection{Evaluación de un software conocido}
	
	Para esta evaluación se selecciona Microsoft Word como ejemplo de software de aplicación.
	
	\begin{longtable}{|p{6cm}|p{7cm}|}
		\hline
		\textbf{Fortalezas} & \textbf{Debilidades} \\
		\hline
		Presenta una interfaz conocida y fácil de usar. & Requiere licencia para utilizar todas sus funciones. \\
		\hline
		Permite elaborar documentos académicos y profesionales. & Puede consumir recursos considerables del sistema. \\
		\hline
		Tiene herramientas de formato, revisión y diseño. & Algunas funciones dependen del ecosistema de Microsoft. \\
		\hline
	\end{longtable}
	
	% ==================================================
	% LENGUAJES
	% ==================================================
	\section{Lenguajes de alto y bajo nivel}
	
	Los lenguajes de programación permiten dar instrucciones a una computadora. Estos pueden clasificarse en lenguajes de alto nivel y lenguajes de bajo nivel.
	
	\subsection{Lenguaje de alto nivel}
	
	Un lenguaje de alto nivel es más cercano al lenguaje humano y facilita la escritura de programas.
	
	\begin{itemize}
		\item Es más fácil de leer y escribir.
		\item Permite desarrollar programas con mayor rapidez.
		\item Reduce la complejidad del código.
	\end{itemize}
	
	\subsection{Lenguaje de bajo nivel}
	
	Un lenguaje de bajo nivel está más relacionado con el hardware y con el lenguaje máquina.
	
	\begin{itemize}
		\item Permite mayor control del hardware.
		\item Puede ofrecer mayor eficiencia.
		\item Requiere mayor conocimiento técnico.
	\end{itemize}
	
	\subsection{Ejemplo en Python}
	
	\begin{verbatim}
		print("Hola mundo")
	\end{verbatim}
	
	\subsection{Comparación con lenguaje ensamblador}
	
	\begin{verbatim}
		MOV AH, 09H
		LEA DX, MENSAJE
		INT 21H
	\end{verbatim}
	
	% ==================================================
	% CONCLUSIONES
	% ==================================================
	\section*{Conclusiones}
	
	El software es un elemento fundamental dentro de los sistemas informáticos, ya que permite realizar tareas, administrar recursos y desarrollar nuevas aplicaciones.
	
	La clasificación del software ayuda a comprender su función dentro de un entorno tecnológico. Además, la calidad del software permite evaluar aspectos importantes como funcionalidad, fiabilidad, usabilidad, eficiencia, mantenibilidad y portabilidad.
	
	Finalmente, los lenguajes de alto y bajo nivel cumplen funciones importantes en la programación. Los lenguajes de alto nivel facilitan el desarrollo, mientras que los lenguajes de bajo nivel permiten mayor control sobre el hardware.
	
	% ==================================================
	% REFERENCIAS APA 7
	% ==================================================
	\clearpage
	
	\ClearShipoutPictureBG
	
	\FondoFinal
	
	\thispagestyle{fancy}
	
	\section*{Referencias}
	
	\begin{hangparas}{1.27cm}{1}
		
		Pressman, R. S., \& Maxim, B. R. (2020). \textit{Ingeniería del software: Un enfoque práctico} (9.a ed.). McGraw-Hill Education.
		
		Sommerville, I. (2016). \textit{Ingeniería de software} (10.a ed.). Pearson Educación.
		
		Stallings, W. (2018). \textit{Sistemas operativos: Aspectos internos y principios de diseño} (9.a ed.). Pearson.
		
		Tanenbaum, A. S., \& Bos, H. (2015). \textit{Sistemas operativos modernos} (4.a ed.). Pearson Educación.
		
		Silberschatz, A., Galvin, P. B., \& Gagne, G. (2018). \textit{Fundamentos de sistemas operativos} (10.a ed.). Wiley.
		
		ISO/IEC. (2011). \textit{ISO/IEC 25010: Systems and software engineering — Systems and software Quality Requirements and Evaluation (SQuaRE)}. International Organization for Standardization.
		
		Oracle Corporation. (2024). \textit{The Java Language Specification}. https://docs.oracle.com/
		
		Massachusetts Institute of Technology. (2023). \textit{Introduction to Computer Science and Programming}. MIT OpenCourseWare. https://ocw.mit.edu/
		
		Harvard University. (2023). \textit{CS50 Introduction to Computer Science}. Harvard OpenCourseWare. https://cs50.harvard.edu/
		
	\end{hangparas}
	
	\ClearShipoutPictureBG
	
\end{document}